\subsection{Bias, correlation shift and support mismatch}
\label{app:bias_correlation}
We first present in \Cref{app:app_bias_full} a decomposition of the \ood bias without any assumptions.
We then justify in \Cref{app:bias_correlation_ass_noiidbias} the simplifying \Cref{ass:no_bias_iid} from \Cref{subsec:expression_ood_bias}.
\subsubsection{\ood bias}
\label{app:app_bias_full}
\begin{proposition}[\ood bias]
    \label{prop:app_bias_full}
    Denoting $\bar{f}_{S}\left(x\right) = \mathbb{E}_{l_S}[f\left(x,\theta\left(l_S\right)\right)]$, the bias is:
    \begin{align*}
        &\mathbb{E}_{(x,y)\sim p_T} [\biasb^2(x, y)] = \int_{\X_T \cap \X_S} \left(f_T\left(x\right)-f_S\left(x\right)\right)^{2} p_T(x) dx                                          & (\text{Correlation shift})                                               \\
                         & + \int_{\X_T \cap \X_S} \left(f_{S}(x)-\bar{f}_{S}\left(x\right)\right)^{2} p_T(x) dx                                                            & (\text{Weighted \iid bias})                                               \\
                         & + \int_{\X_T \cap \X_S} 2 \left(f_T\left(x\right) - f_S\left(x\right)\right) \left(f_S\left(x\right) - \bar{f}_{S}\left(x\right)\right) p_T(x)dx & (\text{Interaction \iid bias and corr. shift})                     \\
                         & + \int_{\X_T \setminus \X_S} \left(f_T\left(x\right)-\bar{f}_{S}\left(x\right)\right)^{2} p_T(x) dx.                                   & (\text{Support mismatch}) %
    \end{align*}%
\end{proposition}%

\begin{proof}
    This proof is original and based on splitting the \ood bias in and out of $\X_S$:
    \begin{align*}
        &\mathbb{E}_{(x,y)\sim p_T}[\biasb^2(x, y)] = \mathbb{E}_{(x,y)\sim p_T} \left(y -\bar{f}_{S}\left(x\right)\right)^{2}                          \\
                                         & = \int_{\X_T} \left(f_T\left(x\right)-\bar{f}_{S}\left(x\right)\right)^{2} p_T(x) dx                \\
                                         & =\int_{\X_T \cap \X_S} \left(f_T\left(x\right)-\bar{f}_{S}\left(x\right)\right)^{2} p_T(x) dx +\int_{\X_T \setminus \X_S} \left(f_T\left(x\right)-\bar{f}_{S}\left(x\right)\right)^{2} p_T(x) dx.
    \end{align*}
    To decompose the first term, we write $\forall x \in \X_S$, $ - \bar{f}_{S}\left(x\right) = - f_S\left(x\right) + \left(f_{S}(x) - \bar{f}_{S}\left(x\right)\right)$.
    \begin{align*}
        &\int_{\X_T \cap \X_S} \left(f_T\left(x\right)-\bar{f}_{S}\left(x\right)\right)^{2} p_T(x) dx=\int_{\X_T \cap \X_S} \left((f_T\left(x\right) - f_S\left(x\right)) + \left(f_{S}(x) - \bar{f}_{S}\left(x\right)\right) \right)^{2} p_T(x) dx \\
        & =\int_{\X_T \cap \X_S} \left(f_T\left(x\right) - f_S\left(x\right)\right)^{2} p_T(x)dx + \int_{\X_T \cap \X_S} \left(f_{S}(x) - \bar{f}_{S}\left(x\right)\right)^{2} p_T(x)dx \\
        & + \int_{\X_T \cap \X_S} 2 \left(f_T\left(x\right) - f_S\left(x\right)\right) \left(f_S\left(x\right) - \bar{f}_{S}\left(x\right)\right) p_T(x)dx.
    \end{align*}%
    % Finally because $\X:=\{x\in \X / p_T(x)>0\}$, we have:
    % $$\int_{\X_T \cap \X_S} \left(f_{S}(x) - \bar{f}_{S}\left(x\right)\right)^{2} p_T(x)dx = \int_{\X_T \cap \X_S} \left(f_{S}(x) - \bar{f}_{S}\left(x\right)\right)^{2} p_T(x)dx.$$
\end{proof}
The four terms can be qualitatively analyzed:
\begin{itemize}
    \item The first term measures differences between train and test labelling function. By rewriting $\forall x \in \X_T \cap \X_S$, $f_T(x) \triangleq \mathbb{E}_{p_T}[Y|X=x]$ and $f_S(x) \triangleq \mathbb{E}_{p_S}[Y|X=x]$, this term measures whether conditional distributions differ.
    This recovers a similar expression to the correlation shift formula from \cite{ye2021odbench}.
    \item The second term is exactly the \iid bias, but weighted by the marginal distribution $p_T(X)$.
    \item The third term  $\int_{\X_T \cap \X_S} 2 \left(f_T\left(x\right) - f_S\left(x\right)\right) \left(f_S\left(x\right) - \bar{f}_{S}\left(x\right)\right) p_T(x)dx$ measures to what extent the \iid bias compensates the correlation shift. It can be negative if (by chance) the \iid bias goes in opposite direction to the correlation shift.
    \item The last term measures support mismatch between test and train marginal distributions. It lead to the \enquote{No free lunch for learning representations for DG} in \cite{ruan2022optimal}. The error is irreducible because \enquote{outside of the source domain, the label distribution is unconstrained}: \enquote{for any domain which gives some probability mass on an example that has not been seen during training, then all \textelp{} labels for that example} are possible.
\end{itemize}

\subsubsection{Discussion of the small \iid bias \Cref{ass:no_bias_iid}}
\label{app:bias_correlation_ass_noiidbias}
\Cref{ass:no_bias_iid} states that $\exists \epsilon > 0 \text{~small~s.t.~} \forall x\in \X_S, |f_{S}\left(x\right)-\bar{f}_{S}\left(x\right)|\leq \epsilon$ where $\bar{f}_{S}\left(x\right) = \mathbb{E}_{l_S} \left[f\left(x,\theta\left(l_S\right)\right)\right]$.
$\bar{f}_{S}$ is the expectation over the possible learning procedures $l_S=\{d_S, c\}$. Thus \Cref{ass:no_bias_iid} involves:
\begin{itemize}
    \item the network architecture $f$ which should be able to fit a given dataset $d_S$. This is realistic when the network is sufficiently parameterized, \ie when the number of weights $|\theta|$ is large.
    \item the expected datasets $d_S$ which should be representative enough of the underlying domain $S$; in particular the dataset size $n_S$ should be large.
    \item the sampled configurations $c$ which should be well chosen: the network should be trained for enough steps, with an adequate learning rate ...
\end{itemize}
For DiWA, this is realistic as it selects the weights with the highest training validation accuracy from each run. For SWAD \cite{cha2021wad}, this is also realistic thanks to their overfit-aware weight selection strategy. In contrast, this assumption may not perfectlty hold for MA \cite{arpit2021ensemble}, which averages weights starting from batch $100$ until the end of training: indeed, $100$ batches are not enough to fit the training dataset.%
\subsubsection{\ood bias when small \iid bias}
\label{app:simplified_bias}
We now develop our equality under \Cref{ass:no_bias_iid}.
\begin{proposition*}[\ref{prop:bias}. \ood bias when small \iid bias]
    With a bounded difference between the labeling functions $f_T-f_S$ on $\X_T \cap \X_S$, under \Cref{ass:no_bias_iid}, the bias on domain $T$ is:
    \begin{equation*} \tag{\ref{eq:bias}}
        \begin{aligned}
            &\mathbb{E}_{(x,y)\sim p_T}[\biasb^2(x, y)] = \text{Correlation shift} + \text{Support mismatch} + O(\epsilon), \\
            &\text{where~} \text{Correlation shift} = \int_{\X_T \cap \X_S} \left(f_T(x)-f_S(x)\right)^{2} p_T(x) dx, \\
            &\text{and~} \text{Support mismatch} = \int_{\X_T \setminus \X_S} \left(f_T(x)-\bar{f}_{S}\left(x\right)\right)^{2} p_T(x) dx.%
        \end{aligned}%
    \end{equation*}%
\end{proposition*}%
\begin{proof}
    We simplify the second and third terms from \Cref{prop:app_bias_full} under \Cref{ass:no_bias_iid}.

    \textbf{The second term} is $\int_{\X_T \cap \X_S} \left(f_{S}\left(x\right)-\bar{f}_{S}\left(x\right)\right)^{2} p_T(x) dx$. Under \Cref{ass:no_bias_iid}, $|f_{S}\left(x\right)-\bar{f}_{S}\left(x\right)|\leq \epsilon$. Thus the second term is $O(\epsilon^2)$.
    
    \textbf{The third term} is $\int_{\X_T \cap \X_S} 2 \left(f_T\left(x\right) - f_S\left(x\right)\right) \left(f_S\left(x\right) - \bar{f}_{S}\left(x\right)\right) p_T(x)dx$.
    As $f_T - f_S$ is bounded on $\X_S \cap \X_T$, $\exists K\geq0$ such that $\forall x \in \X_S$,
    $$
        |\left(f_T(x) - f_S(x)\right)\left(f_{S}(x) - \bar{f}_{S}\left(x\right)\right) p_T(x)| \leq K \left|f_{S}(x) - \bar{f}_{S}\left(x\right)\right|p_T(x) = O(\epsilon) p_T(x).
    $$
    Thus the third term is $O(\epsilon)$.
    
    Finally, note that we cannot say anything about $\bar{f}_{S}\left(x\right)$ when $x\in \X_T \setminus \X_S$.
\end{proof}%

To prove the previous equality, we needed a bounded difference between labeling functions $f_T-f_S$ on $\X_T \cap \X_S$.
We relax this bounded assumption to obtain an inequality in the following \Cref{prop:oodbias_smalliid_nobound}.
\begin{proposition}[\ood bias when small \iid bias without bounded difference between labeling functions]
\label{prop:oodbias_smalliid_nobound}
    Under \Cref{ass:no_bias_iid},
    \begin{equation}
        \begin{aligned}
            &\mathbb{E}_{(x,y)\sim p_T}[\biasb^2(x, y)] \leq 2 \times \text{Correlation shift} + \text{Support mismatch} + O(\epsilon^2)
            % , \\
            % &\text{where~} \text{Correlation shift} = \int_{\X_T \cap \X_S} \left(f_T(x)-f_S(x)\right)^{2} p_T(x) dx, \\
            % &\text{and~} \text{Support mismatch} = \int_{\X_T \setminus \X_S} \left(f_T(x)-\bar{f}_{S}\left(x\right)\right)^{2} p_T(x) dx. %
        \end{aligned}%
    \end{equation}%
\end{proposition}%
\begin{proof}
    % Under \Cref{ass:no_bias_iid}, $|f_{S}\left(x\right)-\bar{f}_{S}\left(x\right)|\leq \epsilon$.
    We follow the same proof as in \Cref{prop:app_bias_full}, except that we now use: $(a+b)^2\leq 2 (a^2 + b^2)$. Then,
    \begin{align*}
    &\int_{\X_T \cap \X_S} \left(f_T\left(x\right)-\bar{f}_{S}\left(x\right)\right)^{2} p_T(x) dx=\int_{\X_T \cap \X_S} \left((f_T\left(x\right) - f_S\left(x\right)) + \left(f_{S}(x) - \bar{f}_{S}\left(x\right)\right) \right)^{2} p_T(x) dx \\
    & \leq 2 \times \int_{\X_T \cap \X_S} \left(f_T\left(x\right) - f_S\left(x\right)\right)^{2} + \left(f_{S}(x) - \bar{f}_{S}\left(x\right)\right)^{2} p_T(x) dx \\
    & \leq 2 \times \int_{\X_T \cap \X_S} \left(f_T\left(x\right) - f_S\left(x\right)\right)^{2} p_T(x) dx + 2 \times \int_{\X_T \cap \X_S} \epsilon^2 p_T(x) dx \\
    & \leq 2 \times \int_{\X_T \cap \X_S} \left(f_T\left(x\right) - f_S\left(x\right)\right)^{2} p_T(x) dx + O(\epsilon^2)
    \end{align*}%
    % Thus the second and third terms from \Cref{prop:app_bias_full} are zero.
    % Then we replace $\forall x \in \X_S \cap  \X_T$: $$f_S(x)=\mathbb{E}_{p_S}[Y|X=x] \text{~~and~~} f_T(x)=\mathbb{E}_{p_T}[Y|X=x].$$
    % Finally, note that we cannot say anything about $\bar{f}_{S}\left(x\right)$ when $x\in \X_T \setminus \X_S$.
\end{proof}%
